\chapter*{How to read this book}
\addcontentsline{toc}{chapter}{How to read this book}
\markboth{How to read this book}{How to read this book}

This book is for the resource geologist who has a block model, a set of drillholes and a question,
and for the reviewer who has to check someone else's answer. It assumes you know what a cutoff, a
swath and a declustered mean are. It does not assume you know micro.

micro opens your files where they are, never uploads them, and keeps a record of what it did with
them. The book follows the same order a piece of work does: look at the data, ask it questions,
put the answers on a page, then make the whole thing repeatable and checkable.

\section*{Paths through it}

\begin{description}[leftmargin=0pt, labelsep=0.6em, font=\sffamily\bfseries\color{ink}]
\item[New to micro.] Read Part~I. \Cref{ch:tour} is a guided tour of the demo project, from opening
  it to a figure and a first routine. It takes about an hour and touches every major idea once.
\item[Here for one task.] Each chapter stands on its own. The index, the glossary
  (\cref{app:glossary}) and the command palette (\key{⌘K} in micro) find the rest.
\item[Reviewing an estimate.] Read Part~IV, \cref{ch:analysis} to \cref{ch:validation}, for what
  each number means and how micro computes it. Then read \cref{ch:checks} and \cref{ch:records} for
  how a routine's checks and run records let you audit someone's month-end without redoing it.
\item[Automating.] Part~VI climbs from one saved analysis (\cref{ch:recipes}) to a month-end
  routine that checks itself, compares with last month and leaves a signed-off record. Developers
  embedding micro in their own page start at \cref{app:host}.
\end{description}

\section*{Follow along}

Every example uses the \term{demo project}. Open it from micro's start screen with \ui{▸ open the
demo project}: a 1.8-million-block iron deposit (\file{deposit.csv}, with \col{FE}, \col{SG} and
\col{LITO}), twelve drillholes drilled in fans across it, the pre-mining topography, three pit phases
and an ore shell. micro generates it on your machine in well under a minute, so you need no data
to work through the book.

The screenshots and charts are made by a script that drives micro on that same demo project, and
the chart values come straight from micro's own tools. When a figure says 24~Mt above 60\%~Fe,
that is what micro computed for this edition. If your copy of micro is newer, a number may differ
slightly. The method won't.

\section*{Conventions}

\begin{center}\small
\begin{tabularx}{\linewidth}{@{}lL@{}}\toprule
\menu{File>Save project as…} & a menu path \\
\ui{Run} & a button, a field, a tab or a checkbox \\
\key{⌘K} & a key. \key{⌘} is \key{Ctrl} on Windows and Linux; the book writes ⌘ throughout \\
\col{FE} & a column \\
\file{deposit.csv} & a file, a folder or a layer named after its file \\
\code{FE > 55 and LITO = "HEMATITE"} & something you type: an expression, a setting, YAML \\
\bottomrule
\end{tabularx}
\end{center}

Notes in the margin\side{like this one} give a cross-reference or a shortcut. Boxes in the text
have one job each:

\begin{tip}
A better way to do what the text just described.
\end{tip}
\begin{note}
A fact worth knowing before you go on.
\end{note}
\begin{caution}
A trap: the place where a number can quietly come out wrong.
\end{caution}
\begin{why}
Why micro works this way. Read these when you want to trust a result, not only reproduce it.
\end{why}
\begin{tryit}
An exercise on the demo project, in steps you can do now. Each starts from the demo as it opens,
or tells you how to get there.
\end{tryit}

\section*{What micro is not}

micro is not an estimation package. It doesn't krige or simulate, and it doesn't design pits. It reads
the outputs of those tools, lets you look at them, measure them and compare them, and writes down what
it did. \Cref{ch:micro} says where its boundaries are and why.
